\tikzset{
  perfent/.style={arqbase, fill=arq-datos, minimum width=34mm, minimum height=6mm,
                  inner sep=2pt, font=\sffamily\scriptsize},
  perfpeso/.style={arqbase, fill=flujo-decision, draw=acc-deficiente, minimum width=9mm,
                   minimum height=5.5mm, inner sep=1pt, font=\sffamily\scriptsize},
  perfmod/.style={arqmodulo, minimum width=9mm, minimum height=5.5mm, inner sep=1pt,
                  font=\sffamily\scriptsize},
  perfclase/.style={arqbase, fill=uoc-azul, text=papel, minimum width=9mm, minimum height=5.5mm,
                    inner sep=1pt, font=\sffamily\scriptsize\bfseries},
  perfcont/.style={arqbase, fill=papel, draw=serie-2, line width=\trazo, rounded corners=2.8mm,
                   minimum width=9mm, minimum height=5.5mm, inner sep=1pt, font=\sffamily\scriptsize},
  perfnota/.style={etiqueta, align=center, font=\sffamily\tiny},
}

% ---------------------------------------------------------------------
% (a) 3.4.1 Diferentes pesos por grupo — Moura et al.
% ---------------------------------------------------------------------
\begin{wrapfigure}{r}{0.36\textwidth}
  \centering
  \vspace{-8pt}
  \begin{tikzpicture}[node distance=2mm]
    \node[perfent] (e) {Indicadores de la calle};
    \node[perfpeso, below=6mm of e] (w2) {$w_2$};
    \node[perfpeso, left=of w2] (w1) {$w_1$};
    \node[perfpeso, right=of w2] (w3) {$w_3$};
    \node[perfclase, below=5mm of w1] (c1) {G$_1$};
    \node[perfclase, below=5mm of w2] (c2) {G$_2$};
    \node[perfclase, below=5mm of w3] (c3) {G$_3$};
    \foreach \i in {1,2,3} {\draw[relacion] (e.south) -- (w\i.north); \draw[relacion] (w\i) -- (c\i);}
    \node[perfnota, right=1mm of w3] {pesos};
    \node[perfnota, right=1mm of c3] {clase};
  \end{tikzpicture}
  \caption{Un índice con pesos $w_i$ por grupo.}
  \label{fig:perfiles-pesos}
  \vspace{-10pt}
\end{wrapfigure}
